\documentclass[12pt]{article}

\usepackage[margin=2cm]{geometry}
\usepackage{natbib}
\usepackage{amsmath, amssymb, bm}
\usepackage{setspace}
\usepackage{xcolor}
\usepackage{longtable}
\usepackage{booktabs}
\usepackage{array}

\newcommand{\E}{\mathbb{E}}

\singlespacing
\setlength{\parindent}{1.5em}
\setlength{\parskip}{0.25em}

\begin{document}

\section{Research Project Background}

\subsection{Motivation and Taiwan Context}

This project studies how external uncertainty shocks originating in the United States, China,
and the global economy are transmitted to Taiwan, which external sources matter most for
Taiwan's domestic uncertainty, and whether the dominant transmission channel changes over
time. These questions are especially important because Taiwan is deeply integrated with both
the U.S. and Chinese economies and lies directly in the path of major policy and geopolitical
developments between them. When the Federal Reserve tightens monetary policy, when
U.S.--China trade disputes escalate, or when cross-strait tensions rise, Taiwan often faces
immediate consequences for trade, investment, capital flows, exchange rates, and financial
conditions.

Major policy shifts and geopolitical realignments rarely move from initial proposal to a
well-defined policy trajectory overnight. Central banks move in sequences of decisions, trade
conflicts unfold through rounds of negotiation and retaliation, and cross-strait relations evolve
over extended periods rather than at a single point in time. As a result, there is often a long
interval between the anticipation of a policy change and the emergence of a stable policy
regime. During this interval, firms, households, and investors face heightened uncertainty about
the future path of policy, demand, and financial conditions, which can affect investment,
consumption, hiring, and portfolio decisions \citep{bloom2009impact,bloom2014fluctuations}.
Because these uncertainty shocks often materialize as high-frequency spikes rather than smooth
quarter-to-quarter transitions, this project uses monthly data to capture the evolving dynamics
without sacrificing the macroeconomic dimension. Monthly observations are granular enough to
capture short-lived uncertainty spikes that quarterly data often smooth out, while still allowing
the analysis to include real-side variables such as industrial production, trade flows, employment,
and prices.

Existing uncertainty measures, such as global financial volatility indices or country-specific
policy uncertainty indices, do not distinguish between the channels through which uncertainty
transmits to the economy. Widely adopted indices such as the VIX or the
\citet{baker2016measuring} Economic Policy Uncertainty (EPU) index compress the
multidimensional nature of uncertainty into a single scalar, obscuring whether a given shock
propagates through real activity or financial markets. For Taiwan, this distinction has direct
policy content:
macroeconomic transmission points to trade flows, export demand, and production linkages,
whereas financial transmission points to capital flows, asset prices, exchange-rate pressure,
and credit conditions.

\subsection{Research Questions}

The project is organized around three main research questions.

\begin{enumerate}
    \item Which external uncertainty sources---U.S., Chinese, or global---contribute most to
    Taiwan's domestic macroeconomic and financial uncertainty?

    \item Do these external shocks transmit to Taiwan mainly through macroeconomic channels
    or through financial channels?

    \item Can the time-varying empirical transmission channels be linked to structural economic
    mechanisms, such as collateral constraints, risk-premium channels, or working-capital
    frictions?
\end{enumerate}

The first question concerns the source of uncertainty. Taiwan is simultaneously exposed to the
United States, China, and global financial and geopolitical conditions. It is therefore not enough
to estimate the effect of a single external index. The empirical model must allow different external
drivers to contribute differently over time.

The second question concerns the channel of transmission. An increase in uncertainty associated
with U.S. monetary policy, for example, may affect Taiwan through weaker export demand and
production expectations, or through exchange-rate depreciation, financial volatility, and tighter
credit conditions. These channels are conceptually distinct and have different policy implications.

The third question links empirical measurement to structural interpretation. A reduced-form
model can identify when external uncertainty affects Taiwan and whether the effect is associated
with macroeconomic or financial conditions. However, it does not by itself reveal the underlying
friction. To address this limitation, the project develops a wedge-based framework that maps
uncertainty shocks into Business Cycle Accounting (BCA) wedges. This framework allows the
project to test whether the data are more consistent with an intertemporal collateral or
risk-premium mechanism, or with a static working-capital mechanism.

\subsection{Empirical and Methodological Challenges}

Identifying the transmission of external uncertainty to Taiwan is difficult for four reasons.

First, external and domestic sources of uncertainty are tightly intertwined, creating a severe
challenge for identification. A single geopolitical or policy event may simultaneously trigger
global financial volatility, shift Taiwan's domestic political environment, and alter expectations
about trade and production. In a small-scale model that lacks sufficient domestic controls, these
domestic movements may be erroneously attributed solely to the external shock. This
omitted-variable bias can lead to a misdiagnosis of the true drivers of uncertainty, overstating
the direct influence of external factors while neglecting local transmission mechanisms
\citep{carriero2018measuring}. Addressing this issue requires a rich information set of both
domestic and external variables.

Second, while a large-scale model is necessary to address this bias, it introduces identification
trade-offs. In large systems, standard Cholesky-based identification becomes problematic because
results are sensitive to variable ordering. Addressing this issue requires an order-invariant
framework. However, order invariance is difficult to reconcile with the strict block-exogeneity
restrictions commonly used in small open economy models. Consequently, the empirical
framework should avoid both arbitrary recursive orderings and rigid zero restrictions that
predetermine how external and domestic variables interact.

Third, many existing approaches face two related limitations. One is their reliance on proxy
measures of uncertainty; the other is their use of a two-step procedure that first estimates
uncertainty and then evaluates its macroeconomic effects in a separate model. Widely used
text-based indices built from international news may not reflect domestic conditions and can
create a perception gap between external coverage and uncertainty as experienced by local
households, firms, and investors. The two-step design also treats the estimated uncertainty
series as observed data, ignores estimation uncertainty, and can induce measurement-error
bias and model inconsistency \citep{carriero2018measuring}. A coherent framework should
therefore estimate uncertainty and its economic effects jointly.


\subsection{Relation to the Literature}

This project is related to three strands of literature. The first studies the macroeconomic effects
of uncertainty shocks. This literature shows that uncertainty can depress investment, hiring, and
economic activity by increasing the option value of waiting and by strengthening precautionary
motives \citep{bloom2009impact,bloom2014fluctuations}. A related strand studies policy
uncertainty and trade policy uncertainty, emphasizing that uncertainty about future policy can
affect firms' investment, export participation, and production decisions even before policy changes
are fully implemented \citep{baker2016measuring,handley2022tpu}.

The second strand develops empirical methods for measuring uncertainty in macroeconomic
models. Existing work shows that large information sets are important because small models can
confound uncertainty shocks with omitted domestic and external conditions
\citep{carriero2018measuring}. The OI-SVMVAR framework of
\citet{davidson2025investigating} is especially relevant because it jointly estimates uncertainty
and its macroeconomic effects, allows uncertainty to enter the conditional mean, and avoids
ordering dependence in large VAR systems.

The third strand is the Business Cycle Accounting literature. BCA decomposes observed aggregate
fluctuations into wedges that summarize distortions in prototype equilibrium conditions
\citep{chari2007business}. This
project uses BCA not as a complete structural model, but as a diagnostic device. The
small-open-economy setting is important because Taiwan's external absorption, imported inputs, and
exchange-rate movements may appear in efficiency and external wedges rather than only in the
standard investment and labor wedges. By deriving how uncertainty-related risk corrections
project into BCA wedges, the project links reduced-form uncertainty measurement to structural
mechanism testing.

Relative to these literatures, the project makes two advances. First, it studies Taiwan as a small
open economy exposed simultaneously to U.S., Chinese, and global uncertainty sources. Second,
it combines time-varying empirical channel classification with wedge-level structural validation.
This combination allows the project to ask not only whether uncertainty matters, but also where
it comes from, how it enters the economy, and which friction best explains its effects.

\subsection{Research Design and Main Contributions}

To address these challenges, the first year of the project applies the order-invariant stochastic
volatility in mean vector autoregression (OI-SVMVAR) framework of
\citet{davidson2025investigating} to Taiwan. The model jointly estimates Taiwan's domestic
macroeconomic uncertainty and financial uncertainty as latent common volatility factors. It also
allows uncertainty to enter the conditional mean of the vector autoregression, so uncertainty can
directly affect macroeconomic and financial outcomes rather than only changing their variances.

The key adaptation to Taiwan is to treat \textbf{external drivers as unclassified variables}. Taiwan's
clearly real-side variables anchor the domestic macroeconomic uncertainty factor, while Taiwan's
clearly financial variables anchor the domestic financial uncertainty factor. External drivers,
including U.S. monetary policy, Chinese activity, global financial risk, geopolitical risk, and
trade-policy uncertainty, are then classified by their time-varying co-movement with Taiwan's
macroeconomic and financial common factors. The resulting posterior classification probabilities
provide a time-varying indicator of the operative transmission channel without requiring the
researcher to impose it ex ante.

This research design differs from a standard country-level application of an existing model. The
unclassified-variable mechanism is used here not only to classify ambiguous domestic variables,
but also to identify how external shocks enter a small open economy. This is the first contribution
of the project.

The second contribution is substantive. The project constructs a monthly Taiwan external
uncertainty dataset covering domestic macroeconomic variables, domestic financial variables,
U.S. indicators, Chinese indicators, global financial and geopolitical risk indicators, and trade-policy
uncertainty measures. This dataset allows the project to compare the relative importance of U.S.,
Chinese, and global uncertainty sources within a single empirical framework.

The third contribution is methodological. The project uses an order-invariant identification
strategy to avoid arbitrary Cholesky ordering in a large model. This feature is important because
the project studies multiple external and domestic variables whose contemporaneous relationships
are difficult to order on purely theoretical grounds.

The fourth contribution is structural. The project develops a wedge-projection framework that
links empirical uncertainty factors to BCA wedges. This framework allows the project to move
beyond reduced-form impulse responses and ask which structural friction is most consistent with
the data. In particular, the project compares two competing mechanisms. A collateral or
risk-premium channel should mainly appear in the investment wedge. A working-capital channel
should mainly appear in the after-tax labor wedge. This wedge-based comparison provides a
testable bridge between the empirical OI-SVMVAR evidence and micro-founded economic
mechanisms.

\subsection{Overview of the Three-Year Project}

The three years of the project are designed as one integrated research program.

The first year focuses on empirical identification. I construct a monthly dataset for Taiwan and
external drivers, estimate the OI-SVMVAR model, extract Taiwan's macroeconomic and financial
uncertainty factors, and compute time-varying classification probabilities for external variables.
The main outputs are Taiwan-specific uncertainty factors, posterior classification probabilities,
impulse responses (IRFs), and forecast error variance decompositions (FEVDs).

The second year focuses on theoretical mapping. I develop a wedge-projection framework that
shows how second-moment uncertainty corrections in a nonlinear structural model can appear as
first-order wedges in a BCA prototype. This theoretical component derives wedge signatures for
competing mechanisms. In particular, it distinguishes an intertemporal collateral or risk-premium
channel from a static working-capital channel.

The third year focuses on empirical validation. I construct BCA wedges for Taiwan and estimate
the responses of these wedges to the uncertainty factors extracted in the first year. The estimation
passes the posterior uncertainty from the first-year OI-SVMVAR into the third-year local
projections, rather than using only a single posterior mean path. This design reduces the risk that
the final mechanism comparison is driven by generated-regressor bias.

Taken together, the project provides a unified account of external uncertainty transmission to
Taiwan. The first year identifies when and through which channel external uncertainty enters the
domestic economy. The second year derives structural wedge signatures for competing mechanisms.
The third year tests these signatures in Taiwan's data and translates the results into policy-relevant
evidence.

\section{Preliminary Progress and Feasibility}

\subsection{Preliminary Progress}

The project is feasible because it builds on substantial preliminary work. The first-year dataset
has moved beyond the planning stage. I have completed an automated raw-data collection
pipeline for nearly all variables in the first-year OI-SVMVAR design. The current working design
uses a baseline monthly panel covering 2003M1--2025M12, with a longer price-index robustness
sample beginning in 2001M1. The baseline contains 52 variables: 22 Taiwan macroeconomic
anchors, 8 Taiwan financial anchors, and 22 unclassified variables that include external drivers
and Taiwan-domestic variables whose channel may vary over time. The data appendix documents
the source groups, transformation rules, treatment of upstream price data, and reproducibility
pipeline.

The raw data collection already covers Taiwan macroeconomic indicators from the
Directorate-General of Budget, Accounting and Statistics (DGBAS), Ministry of Economic
Affairs (MOEA), Ministry of Finance (MOF), and National Development Council (NDC) web
sources; Taiwan financial and credit variables from the Central Bank of the Republic of China
(CBC), the Taiwan Stock Exchange (TWSE), Federal Reserve Economic Data (FRED), and
housing-price sources; U.S. variables from FRED and PolicyUncertainty.com; China variables
from National Bureau of Statistics (NBS) and People's Bank of China (PBoC) based sources; and
global variables from FRED, Bank for International Settlements (BIS), the Iacoviello geopolitical
risk dataset, and PolicyUncertainty.com. The remaining first-year data task is therefore not basic data collection,
but construction of the final estimation matrix: harmonizing dates, applying transformations,
materializing derived variables such as the term spread and exchange-rate volatility, documenting
short endpoint gaps, and producing the balanced panel used in the Bayesian estimation.

The most important preliminary progress is on the second-year theoretical component. I have
already developed a standalone wedge-projection report that turns the informal idea that
``uncertainty appears as a BCA wedge'' into a formal theoretical framework. Because
intertemporal financial frictions primarily distort the household Euler equation, their natural
BCA counterpart is the investment wedge. The central result is a projection equation for the BCA
investment wedge:
\[
\hat{\tau}_{x,t}
=
\omega_\tau \mathrm{E}_t \hat{\tau}_{x,t+1}
+
\frac{1}{2}\Phi(\Omega_t),
\]
where $\Phi(\Omega_t)$ is the second-order risk correction generated by stochastic volatility in
the true nonlinear structural model, $\omega_\tau$ is a forward-looking coefficient related to
the steady-state discount factor, capital depreciation, and adjustment-cost curvature, and
$\Omega_t$ denotes the stochastic-volatility state. Under an AR(1) law of motion for volatility,
this equation admits a closed-form solution. This result shows that a second-moment uncertainty
signal in a nonlinear Dynamic Stochastic General Equilibrium (DSGE) model can appear as a
persistent first-order investment wedge in a BCA prototype.

The preliminary theoretical work also derives distinct wedge signatures for two competing
mechanisms. In a Fern\'andez-Villaverde--Guerr\'on-Quintana-style collateral or risk-premium
channel, uncertainty affects the intertemporal Euler equation through second-order risk
corrections \citep{fernandez-villaverde2011risk,fernandez-villaverde2020uncertainty}. This
mechanism projects mainly into the investment wedge. In contrast, a working-capital channel affects firms' static
labor-demand condition because firms must finance part of their wage bill before production.
This mechanism projects mainly into the after-tax labor wedge. These two mechanisms therefore
generate different empirical predictions.

The preliminary work further extends the projection argument to third order. This extension is
important because uncertainty shocks enter higher-order perturbation solutions through
volatility innovations and state-dependent terms. The third-order extension introduces
vol-of-vol effects, level-volatility interaction terms, and asymmetric responses to positive and
negative uncertainty innovations. These objects provide additional empirical handles for
distinguishing structural mechanisms.

The current numerical validation also clarifies the remaining tasks. Preliminary calculations show
that the theoretical closed-form impact object should not be compared mechanically with the
peak of a long-horizon local projection. The next revision will compare the full horizon profile,
including the impact response, sign changes, and peak response. This issue does not weaken the
project. Instead, it provides a clear and manageable task for the second-year theoretical and
numerical validation stage.

\subsection{Feasibility of the Proposed Plan}

The proposed three-year plan is feasible for four reasons.

First, the first-year empirical model is based on an existing and published econometric framework.
The project does not require developing a new Bayesian sampler from scratch. The main task is
to adapt the OI-SVMVAR framework to the Taiwan setting, construct the dataset, and reinterpret
the unclassified-variable mechanism as a transmission-channel identifier for a small open economy.

Second, the data construction is manageable because the main raw-data downloads have already
been implemented. The repository now contains source-specific scripts for FRED, CBC, TWSE,
DGBAS, the Taiwan statistical portal, BIS, China indicators, and CSV-based uncertainty measures.
These scripts have produced the raw series needed for the 52-variable design. The remaining data
risks are limited and well-defined: foreign institutional net buying begins after the baseline
start, global EPU currently has a short endpoint gap, and the China block requires continued
source verification. These issues can be handled through transparent sample trimming,
robustness specifications, or reduced China-block benchmarks without changing the research
design.

Third, the second-year theoretical component already has a clear foundation. The preliminary
report has derived the main projection equation, the closed-form solution under AR(1) volatility,
the distinction between investment-wedge and labor-wedge signatures, and the need for a
third-order extension. The remaining work is to refine the proofs, standardize notation, improve
the numerical validation, and adapt the theoretical presentation to the Taiwan application.

Fourth, the third-year empirical validation uses outputs from the first two years rather than
requiring a separate data project. Once the first-year OI-SVMVAR produces posterior draws of
the uncertainty factors and the second-year theory produces wedge signatures, the third-year task
is to construct Taiwan's BCA wedges and estimate wedge-level local projections. The design also
passes first-year posterior uncertainty into the third-year local projections, which reduces the
risk of generated-regressor bias.

Overall, the project has a clear fallback structure. If the full 52-variable OI-SVMVAR is slow to
converge, I will first estimate a reduced benchmark model using the core Taiwan anchors and the
most important external drivers. If some interpolated variables affect the results, I will report
robustness checks that exclude them. If the empirical wedge responses do not sharply select one
mechanism, I will report the relative strength of evidence across mechanisms rather than forcing
a binary conclusion. These safeguards make the project ambitious but tractable.

\section{The First Year: Methods, Procedures, and Implementation Schedule}

\subsection{Purpose of the First Year}

The first year identifies how external uncertainty shocks enter Taiwan and whether they operate
through macroeconomic or financial channels. The empirical task is not simply to construct an
uncertainty index. Instead, the goal is to jointly estimate Taiwan's domestic uncertainty factors,
measure the effects of these factors on macroeconomic and financial variables, and classify the
transmission channel of external drivers over time.

The first-year analysis has four steps. First, I construct a monthly dataset that combines Taiwan's
domestic macroeconomic variables, domestic financial variables, U.S. indicators, Chinese
indicators, global risk measures, and trade-policy uncertainty measures. Second, I estimate an
order-invariant stochastic volatility in mean vector autoregression (OI-SVMVAR). Third, I extract
two latent domestic uncertainty factors: Taiwan macroeconomic uncertainty and Taiwan financial
uncertainty. Fourth, I use the posterior classification probabilities of unclassified variables to
identify whether each external driver is transmitted to Taiwan through the macroeconomic or
financial channel at each point in time.

This first-year component produces the empirical foundation for the whole project. The extracted
uncertainty factors become the key inputs for the third-year BCA validation. The time-varying
classification probabilities also provide the first direct evidence on whether U.S., Chinese, and
global shocks affect Taiwan through different channels in different episodes.

\subsection{Empirical Strategy Overview}

The empirical framework follows the OI-SVMVAR model of
\citet{davidson2025investigating}. Let $y_t$ denote the $N \times 1$ vector of
monthly observables. In the baseline application, $N=52$. The observation equation is
\[
y_t
=
\sum_{\ell=1}^{p} B_\ell y_{t-\ell}
+
\sum_{j=0}^{q} A^h_j h_{t-j}
+
B_0^{-1}\varepsilon^y_t,
\qquad
\varepsilon^y_t \sim N(0,U_t),
\]
where
\[
h_t = (h_{m,t},h_{f,t})'
\]
contains Taiwan's macroeconomic uncertainty factor $h_{m,t}$ and financial uncertainty factor
$h_{f,t}$. The stochastic-volatility-in-mean term,
\[
\sum_{j=0}^{q} A^h_j h_{t-j},
\]
allows uncertainty to affect the conditional mean of the economy. This feature is essential for
the project because uncertainty may change investment, employment, asset prices, credit
conditions, and exchange rates, not merely the variance of these variables.
Appendix~\ref{app:oi-svmvar} provides the compact model statement, including the
log-volatility decomposition, the latent classification states, and the posterior sampling objects.

The contemporaneous impact matrix $B_0$ is unrestricted apart from standard normalizations.
This avoids the arbitrary ordering problem that arises in Cholesky-based VARs. In a large model
with many domestic and external variables, no single recursive ordering has a strong economic
justification. The order-invariant structure is therefore important for the credibility of the impulse
responses, forecast error variance decompositions, and classification probabilities.

The conditional variance matrix $U_t$ is diagonal and is organized into three blocks:
\[
U_t
=
\begin{pmatrix}
\Omega_{m,t} & 0 & 0 \\
0 & \Omega_{u,t} & 0 \\
0 & 0 & \Omega_{f,t}
\end{pmatrix}.
\]
The macroeconomic block anchors $h_{m,t}$, the financial block anchors $h_{f,t}$, and the
unclassified block is allowed to switch between the two factors. For a pre-classified
macroeconomic variable, the common volatility component is $h_{m,t}$. For a pre-classified
financial variable, the common volatility component is $h_{f,t}$. For an unclassified variable,
the common volatility component depends on a latent state:
\[
\omega^u_{i,t}
=
\eta^u_{i,t}
+
h_{s_{i,t},t},
\qquad
s_{i,t}\in\{m,f\}.
\]
The state $s_{i,t}$ follows a two-state Markov chain. The key object for this project is the
posterior classification probability
\[
\pi_{i,t}
=
P(s_{i,t}=m \mid y_{1:T}).
\]
The idiosyncratic term $\eta^u_{i,t}$ captures variable-specific volatility that is not absorbed by
the common macroeconomic or financial uncertainty factor.
When $\pi_{i,t}$ is close to one, the model assigns unclassified variable $i$ to the macroeconomic
uncertainty factor at time $t$. When $\pi_{i,t}$ is close to zero, the model assigns it to the
financial uncertainty factor. For an external variable, this probability is interpreted as a
time-varying indicator of whether the variable is more closely associated with Taiwan's
macroeconomic or financial uncertainty factor. The classification probability is therefore a
channel-association measure; the direction and quantitative importance of external shocks are
evaluated jointly with forecast error variance decompositions and impulse responses.

A natural alternative would be to impose a small-open-economy block-exogeneity restriction,
under which external variables can affect Taiwan contemporaneously but Taiwan variables cannot
affect external variables. I do not impose this restriction in the baseline model for two reasons.
First, block exogeneity would introduce new zero restrictions into the contemporaneous impact
matrix and would compromise the order-invariant structure that motivates the OI-SVMVAR.
Second, Taiwan's role in global semiconductor supply chains makes strict contemporaneous
zero restrictions less innocuous than in a textbook small open economy. The project instead treats
the small-open-economy assumption as an empirical diagnostic. If Taiwan's uncertainty factors
explain little of the forecast error variance of external variables, the small-open-economy
restriction is approximately supported by the data. If they explain a non-negligible share, imposing
the restriction ex ante would have been misleading.

\subsection{Variable Classification and Taiwan Adaptation}

The main innovation in the first-year empirical design is the way the unclassified-variable
mechanism is used in a small open economy. In the original OI-SVMVAR framework, unclassified
variables are domestic variables whose macroeconomic or financial nature is ambiguous. In this
project, the same statistical device is repurposed to study external transmission. External drivers
are placed in the unclassified block because their transmission channel is precisely the object of
interest.

Taiwan's clearly real-side variables anchor the macroeconomic uncertainty factor. These include
industrial production, manufacturing production, export orders, retail sales, food-service sales,
exports, imports, price indices, unemployment, employment, and real wages. Taiwan's clearly
financial variables anchor the financial uncertainty factor. These include the overnight interbank
rate, the ten-year government bond yield, the term spread, stock-market returns, realized equity
volatility, turnover, foreign institutional net buying, and margin buying.

The unclassified block contains three types of variables. The first type is external drivers, including
U.S. monetary and real indicators, Chinese real and monetary indicators, global financial and
geopolitical risk measures, and trade-policy uncertainty measures. The second type is Taiwan
domestic variables whose interpretation can vary over time, such as monetary aggregates, credit
aggregates, the exchange rate, the real effective exchange rate, the TAIEX level, housing prices,
and short-selling balances. The third type is policy and market variables that can behave as either
macroeconomic or financial indicators depending on the episode.

This classification is central to the research design. Pre-classifying the U.S. Federal Funds Rate
as financial, or China's industrial production as macroeconomic, would impose the answer before
estimation. In reality, U.S. monetary tightening may affect Taiwan through financial channels such
as capital flows and exchange-rate pressure, but it may also affect Taiwan through macroeconomic
channels such as export demand and production expectations. Similarly, a slowdown in China may
operate through Taiwan's trade and production networks in one episode, and through asset prices
and risk sentiment in another. By placing these variables in the unclassified block, the model lets
the data determine the operative channel.

The resulting classification probabilities provide a time-varying narrative of external transmission.
For example, if the posterior probability for the U.S. Federal Funds Rate is close to one during a
given period, the model indicates that U.S. monetary uncertainty is more closely associated with
Taiwan's macroeconomic uncertainty. If the probability is close to zero, the model indicates that
the same external driver is more closely associated with Taiwan's financial uncertainty. Plotting
these probabilities across episodes such as the Global Financial Crisis, the U.S.--China trade war,
the COVID-19 shock, and the 2022 Federal Reserve tightening cycle allows the project to compare
how Taiwan's exposure changes over time.

\subsection{Data Construction and Transformations}

The baseline dataset is a balanced monthly panel covering 2003M1--2025M12. The current design
contains 52 variables: 22 macroeconomic variables, 8 financial variables, and 22 unclassified
variables. A longer price-index robustness sample covers 2001M1--2025M12 for specifications
that do not require the TWSE total-return index. The baseline starts in 2003 because the TWSE
total-return index, which provides the preferred equity-return measure, begins in January 2003.
This starting point also captures Taiwan's post-WTO integration period while avoiding the earlier
Asian financial crisis as a major structural break.

The monthly frequency is central to the empirical design. Quarterly data would be too coarse
to capture short uncertainty episodes associated with policy announcements, tariff negotiations,
geopolitical events, and financial-market stress. Daily data would capture financial volatility but
would exclude many real-side variables needed to anchor the macroeconomic uncertainty factor.
Monthly data therefore provide the appropriate compromise: they preserve enough time-series
variation to identify short-lived uncertainty episodes while still allowing the model to include
production, trade, employment, prices, and credit variables.

\begin{center}
\begin{tabular}{p{0.18\textwidth}p{0.42\textwidth}p{0.28\textwidth}}
\toprule
Block & Main variables & Main sources \\
\midrule
Taiwan macroeconomic block &
Industrial production, manufacturing production, export orders, retail sales, exports, imports,
business-cycle indicators, prices, unemployment, employment, real wages, working hours, and
residential construction permits &
DGBAS, MOEA, MOF, NDC \\
\addlinespace
Taiwan financial block &
Overnight interbank rate, 10-year government bond yield, term spread, TAIEX returns,
TAIEX realized volatility, turnover, foreign institutional net buying, and margin buying &
CBC, TWSE \\
\addlinespace
Unclassified Taiwan variables &
CBC discount rate, money aggregates, bank loans, consumer loans, TAIEX level, housing prices,
exchange rates, exchange-rate volatility, real effective exchange rate, and short-selling balance &
CBC, TWSE, FRED, BIS, housing-price sources \\
\addlinespace
External variables &
U.S. Federal Funds Rate, U.S. industrial production, U.S. credit spread, U.S. EPU, China
industrial production, China PPI, China M2, VIX, global geopolitical risk, global EPU, and
U.S. trade-policy uncertainty &
FRED, NBS/PBoC-based sources, PolicyUncertainty.com, GPR dataset, BIS \\
\bottomrule
\end{tabular}
\end{center}

The Taiwan macroeconomic block has been downloaded from the DGBAS nstatdb system and the
official Taiwan statistical portal, supplemented by MOEA, MOF, and NDC series. It includes
production, trade, prices, labor-market, wage, working-hour, business-cycle, and housing-permit
variables. The financial block uses CBC and TWSE data, including interest rates, monetary and
credit aggregates, TAIEX total-return and price-index measures, turnover, foreign investor flows,
margin buying, and short-selling balances. The external block uses FRED, NBS/PBoC-based China
series, BIS, PolicyUncertainty.com, and the Iacoviello geopolitical risk dataset. This block includes
U.S. monetary, real, credit, EPU, and trade-policy uncertainty variables; China industrial
production, PPI, and M2; VIX; global EPU; global geopolitical risk; and Taiwan's real effective
exchange rate.

All series are transformed before estimation to improve stationarity and comparability. Real-side
level variables, price indices, employment, wages, money, and credit variables are expressed as
year-on-year percentage changes. Interest-rate levels are differenced when they display persistent
nonstationarity, while spreads are retained in levels. Equity returns are monthly log returns.
Realized volatility measures are computed from daily returns within each month. Exchange-rate
and equity-market levels are retained as log levels only when their level is economically meaningful
for the classification exercise. All transformed variables are standardized to have zero mean and
unit variance before estimation.

The data appendix documents the source of each series, the transformation rule, and the
reproducibility pipeline. It also records special data-construction choices. For example, Taiwan's
legacy wholesale price index (WPI) ends in 2022M12, while the newer producer price index
(PPI) begins in 2021M1. Because the two series overlap from 2021M1 to 2022M12, I have
constructed a spliced upstream price-pressure series and retained overlap diagnostics. This series
is labeled as an auxiliary upstream price measure rather than as an official continuation of the WPI.

At the current stage, the remaining data issues are narrow rather than fundamental. Foreign
institutional net buying begins in 2004M4, so the baseline will either trim this series, treat it as
a robustness variable, or use a reduced balanced sample when it is included. Global EPU is
currently missing the final month of 2025, which can be handled by waiting for the source update
or ending the common sample in 2025M11. The term spread is fully determined by the downloaded
10-year government bond yield and overnight interbank rate and will be materialized in the final
dataset-building script. For the China block, the current raw-data pipeline has already obtained
industrial production, PPI, and M2 series; the remaining task is source validation and robustness
documentation.

\subsection{Estimation, Inference, and Main Outputs}

The model is estimated using the Bayesian Markov Chain Monte Carlo (MCMC) algorithm
developed for the OI-SVMVAR framework. The estimation jointly samples the VAR coefficients, the contemporaneous impact
matrix, the common log-volatility factors, idiosyncratic volatility components, and the Markov
classification paths for unclassified variables. I will use the implementation details of
\citet{davidson2025investigating} as the baseline and adapt the code to the Taiwan dataset.
Because the baseline system contains 52 variables, the estimation will rely on the shrinkage-prior
structure of the OI-SVMVAR framework to regularize the high-dimensional VAR coefficients and
limit overfitting while preserving the large information set needed to reduce omitted-variable
bias.
For the unclassified variables, the latent macro-financial classification indicators will be treated
as finite-state Markov paths and sampled with a forward-filtering/backward-sampling step, so that
classification probabilities reflect the full posterior path uncertainty rather than a sequence of
period-by-period labels.

The main posterior objects are the following:
\[
\{h_{m,t},h_{f,t}\}_{t=1}^{T},
\qquad
\{\pi_{i,t}\}_{i\in \mathcal{U},t=1}^{T},
\]
where $\mathcal{U}$ denotes the set of unclassified variables. The first object gives Taiwan's
macroeconomic and financial uncertainty factors. The second object gives the time-varying
classification probabilities for external and ambiguous domestic variables.

The first-year analysis will report four sets of results. First, I will plot Taiwan's macroeconomic
and financial uncertainty factors and compare their movements across major episodes. Second, I
will plot classification probabilities for key external variables, including U.S. monetary indicators,
Chinese real indicators, global financial risk, geopolitical risk, and trade-policy uncertainty. Third,
I will compute impulse responses to macroeconomic and financial uncertainty shocks. Fourth, I
will compute forecast error variance decompositions to quantify how much U.S., Chinese, and
global uncertainty shocks contribute to Taiwan's domestic uncertainty factors and observed
economic variables.

The analysis will also include robustness checks. I will compare the full 52-variable model with a
reduced benchmark specification that keeps the core Taiwan anchors and the most important
external drivers. I will check alternative lag lengths, alternative sample starts, and specifications
that exclude interpolated variables. For equity-market variables, I will compare the baseline
TWSE total-return index with a longer price-index-based robustness sample when appropriate.

\subsection{First-Year Work Plan and Anticipated Results}

The first year is organized into four stages.

\begin{center}
\begin{tabular}{p{0.18\textwidth}p{0.36\textwidth}p{0.34\textwidth}}
\toprule
Period & Main task & Expected output \\
\midrule
Months 1--4 &
Finalize the estimation dataset from the completed raw-data pipeline &
Balanced panel, codebook, transformation log, stationarity checks, endpoint-gap decisions \\
\addlinespace
Months 3--7 &
Adapt and test the OI-SVMVAR estimation code &
Validated MCMC workflow and convergence diagnostics \\
\addlinespace
Months 5--9 &
Estimate benchmark and full models &
Posterior draws of $h_{m,t}$, $h_{f,t}$, and classification paths \\
\addlinespace
Months 9--12 &
Analyze transmission channels and uncertainty effects &
Classification plots, impulse responses, variance decompositions, and
conference-paper draft \\
\bottomrule
\end{tabular}
\end{center}

By the end of the first year, I expect to complete four outputs. First, I will produce a documented
monthly dataset for Taiwan and external uncertainty drivers. Second, I will estimate Taiwan's
macroeconomic and financial uncertainty factors using the OI-SVMVAR framework. Third, I will
produce time-varying classification probabilities for U.S., Chinese, global, and trade-policy
variables. Fourth, I will prepare a first empirical paper draft or conference paper that documents
the main first-year findings.

The first-year results will directly feed into the third year. The posterior draws of
$h_{m,t}$ and $h_{f,t}$ will be passed into the BCA wedge validation exercise, so the project can
evaluate structural mechanisms while accounting for the uncertainty generated in the first-stage
estimation.

\section{The Second Year: Methods, Procedures, and Implementation Schedule}

\subsection{Purpose of the Second Year}

The first-year OI-SVMVAR identifies Taiwan's macroeconomic and financial uncertainty factors
and shows whether external drivers transmit mainly through macroeconomic or financial channels.
However, this empirical evidence alone does not identify the underlying structural friction. A
financial-channel response, for example, may reflect a collateral constraint, a time-varying risk
premium, a working-capital requirement, or other financial frictions. The purpose of the second
year is to build a theoretical bridge between the empirical uncertainty factors and structural
economic mechanisms.

The second-year component develops a wedge-projection framework. The framework asks how
second-moment uncertainty terms in a nonlinear structural model appear to an econometrician
who observes the data through a first-order Business Cycle Accounting (BCA) prototype. The key
idea is that uncertainty shocks may not appear as conventional first-order shocks in the true
structural model, but their risk corrections can be projected into first-order BCA wedges. These
wedges then provide diagnostic signatures for different transmission mechanisms.

The baseline nonlinear environment follows the small-open-economy framework with stochastic
volatility developed by Fern\'andez-Villaverde, Guerr\'on-Quintana, and co-authors
\citep{fernandez-villaverde2011risk,fernandez-villaverde2020uncertainty}. In that environment,
uncertainty affects real allocations through intertemporal collateral and risk-premium terms:
tighter collateral risk, precautionary behavior, and time-varying risk premia distort the Euler
equation and should therefore be visible primarily as an investment-wedge signature in a
first-order BCA representation. I will contrast this baseline with a working-capital
counter-hypothesis in which firms finance wage payments before production, so that higher
external finance costs operate mainly through the static labor-demand margin and appear as an
after-tax labor-wedge deterioration rather than an investment-wedge response.

This step is essential for the three-year project. Without it, the first-year evidence would identify
when external uncertainty affects Taiwan and through which reduced-form channel, but it would
not tell us which friction drives the response. The second year therefore supplies the theoretical
mapping needed for the third-year empirical validation.

\subsection{Second-Moment Wedge Projection}

The central theoretical object is the BCA investment wedge. In a first-order BCA prototype, the
investment wedge absorbs distortions that prevent the representative-household Euler equation
from matching the observed path of consumption, investment, and capital returns. If the true
data-generating process is a nonlinear structural model with stochastic volatility, the true Euler
equation contains second-order risk corrections. The BCA prototype, which is first order, cannot
represent these risk corrections directly. Instead, it projects them into wedges. Appendix~\ref{app:wedge-projection}
states the economic model behind this projection, including the BCA Euler equation, the
second-order risk correction, the closed-form projection, and the working-capital counter-hypothesis.

The preliminary theoretical work has derived the following projection equation:
\[
\hat{\tau}_{x,t}
=
\omega_\tau \mathrm{E}_t\hat{\tau}_{x,t+1}
+
\frac{1}{2}\Phi(\Omega_t),
\]
where $\hat{\tau}_{x,t}$ is the log-deviation of the BCA investment wedge,
$\omega_\tau$ is a steady-state coefficient, $\Omega_t$ denotes the stochastic-volatility state,
and $\Phi(\Omega_t)$ is the second-order risk correction in the true structural Euler equation.

This equation has a simple interpretation. A first-order BCA observer does not see the nonlinear
risk correction directly. Instead, the observer recovers a forward-looking investment wedge that
moves with the volatility state. If the volatility state follows an AR(1) process, the projection has
a closed-form solution:
\[
\hat{\tau}_{x,t}
=
\frac{\Phi(\bar{\Omega})}{2(1-\omega_\tau)}
+
\frac{\phi}{2(1-\omega_\tau\rho_\Omega)}
(\Omega_t-\bar{\Omega}),
\]
where $\phi=\Phi'(\bar{\Omega})$ and $\rho_\Omega$ is the persistence of the volatility state.
This expression shows that persistent volatility can create a persistent first-order investment
wedge even when the underlying mechanism operates through second-moment risk corrections.

The second-year task is to formalize this projection result, standardize the notation, and adapt
the theorem to the small-open-economy setting relevant for Taiwan. The theorem provides the
conceptual foundation for linking the first-year uncertainty factors to third-year BCA wedges.

\subsection{Competing Mechanisms and Wedge Signatures}

The wedge-projection framework is useful because different structural frictions imply different
wedge signatures. The second year focuses on two benchmark mechanisms: an intertemporal
collateral or risk-premium channel and a static working-capital channel. It also treats efficiency
and external wedges as Taiwan-specific diagnostic objects because imported intermediate inputs,
trade flows, and exchange-rate movements are central to a small open economy.

The first mechanism is a collateral or risk-premium channel. In this environment, higher financial
uncertainty affects the intertemporal Euler equation through second-order risk corrections. These
corrections may arise from precautionary motives, time-varying risk premia, collateral-risk terms,
or aggregation effects in a two-agent economy. Because the distortion enters the intertemporal
condition, the BCA projection appears mainly in the investment wedge. Under the baseline
collateral-risk mechanism, financial uncertainty is expected to raise the investment wedge.

The second mechanism is a working-capital channel. In this environment, firms must finance part
of their wage bill before production. A rise in external finance spreads then acts like a tax on labor
demand. This mechanism does not primarily distort the intertemporal Euler equation. Instead, it
distorts the static labor-demand condition. Therefore, the BCA projection appears mainly in the
after-tax labor wedge. When financing spreads rise, the after-tax labor wedge should deteriorate.
For Taiwan, however, working capital may also finance imported intermediate inputs rather than
only wage payments. In that case, the same financing friction may also appear in the efficiency
wedge or the government/external wedge, because the distortion affects materials, imported
inputs, and external absorption. The empirical test will therefore treat the labor wedge as the
benchmark signature of a wage-finance channel while using the efficiency and external wedges
as diagnostic objects for input-finance and trade-related variants of working-capital frictions.

The main wedge signatures are summarized as follows:
\begin{center}
\begin{tabular}{p{0.25\textwidth}p{0.32\textwidth}p{0.28\textwidth}}
\toprule
Mechanism & Main economic friction & Predicted BCA signature \\
\midrule
Collateral or risk-premium channel &
Intertemporal Euler distortion from uncertainty risk corrections &
Investment wedge response \\
\addlinespace
Working-capital channel &
Static labor-demand distortion from financing costs &
After-tax labor wedge response \\
\addlinespace
SOE input-finance or trade variant &
Imported-input finance, exchange-rate movements, or external absorption &
Efficiency or external-wedge response \\
\bottomrule
\end{tabular}
\end{center}

This table is the theoretical bridge between the first-year OI-SVMVAR and the third-year BCA
validation. If the empirical uncertainty factors mainly move the investment wedge, the evidence
supports an intertemporal collateral or risk-premium mechanism. If they mainly move the
after-tax labor wedge, the evidence supports a working-capital mechanism. If they mainly move
the efficiency or external wedge, the evidence suggests that Taiwan's small-open-economy
structure, imported-input dependence, or external absorption margin is central to the transmission
process. If several wedges respond, the data may indicate that several frictions operate at the
same time.

\subsection{Third-Order Extension}

The second-year work also extends the projection framework to third order. This extension is
needed because many stochastic-volatility DSGE models use third-order perturbation to capture
the independent role of uncertainty shocks. At second order, the volatility state affects risk
corrections, but the innovation to volatility may not generate the full impact response observed
in generalized impulse responses or local projections. Third-order terms introduce additional
objects that are empirically relevant.

The preliminary third-order extension introduces three new components. The first is a
vol-of-vol correction, which refines the linear coefficient on the volatility state. The second is a
state-dependent level-volatility interaction term, which allows the effect of uncertainty to depend
on the current state of the economy. The third is an asymmetry term, which allows positive and
negative volatility innovations to have responses that are not mirror images of one another.

These objects are useful for empirical interpretation. If the investment-wedge response is stronger
during periods of tight credit conditions, this pattern would be consistent with a state-dependent
collateral mechanism. If positive uncertainty shocks generate larger responses than negative
uncertainty shocks, this asymmetry would support a nonlinear volatility mechanism. These
predictions give the third-year empirical analysis more diagnostic power than a simple comparison
of average impulse responses.

The second-year task is not to make the theoretical model unnecessarily complex. Instead, it is to
derive a disciplined hierarchy of tests. The second-order projection provides the baseline wedge
signature. The third-order extension provides additional checks for state dependence, asymmetry,
and volatility-of-volatility effects.

\subsection{Numerical Validation and Remaining Tasks}

The preliminary numerical work has already identified an important issue. The closed-form
projection coefficient is a local object. It should be compared with the impact response or the
full horizon profile of the empirical wedge response, not mechanically with the maximum response
at a long horizon. A peak response from a local projection may occur many months after the
initial shock and may reflect persistence, propagation, and other state dynamics. Therefore, the
second-year validation will compare the theoretical and empirical objects more carefully.

The numerical validation will proceed in four steps. First, I will compute the theoretical
second-order wedge coefficients under the benchmark collateral and working-capital mechanisms.
Second, I will simulate the structural models and recover the implied BCA wedges from the
simulated data. Third, I will compare theoretical signs with simulated wedge responses at impact,
at the zero crossing, and at the peak horizon. Fourth, I will evaluate the third-order predictions
by studying state-dependent responses and asymmetric responses to positive and negative
uncertainty innovations.

This validation strategy turns the theory into a practical empirical tool. It also reduces the risk
that the third-year mechanism comparison relies on an incorrect theoretical object.

\subsection{Second-Year Work Plan and Anticipated Results}

The second year is organized into four stages.

\begin{center}
\begin{tabular}{p{0.18\textwidth}p{0.36\textwidth}p{0.34\textwidth}}
\toprule
Period & Main task & Expected output \\
\midrule
Months 1--3 &
Finalize the second-moment wedge-projection theorem &
Formal theorem, proof, and notation-consistent derivation \\
\addlinespace
Months 3--6 &
Derive wedge signatures for benchmark mechanisms &
Investment-wedge and labor-wedge diagnostic table \\
\addlinespace
Months 6--9 &
Develop and refine the third-order extension &
State-dependence, asymmetry, and vol-of-vol predictions \\
\addlinespace
Months 9--12 &
Conduct numerical validation and prepare theory paper &
Validated theoretical appendix and standalone theory manuscript \\
\bottomrule
\end{tabular}
\end{center}

By the end of the second year, I expect to complete four outputs. First, I will produce a formal
wedge-projection theorem that maps stochastic-volatility risk corrections into BCA wedges.
Second, I will derive a wedge-signature table that distinguishes collateral, risk-premium, and
working-capital mechanisms. Third, I will complete numerical validation using simulated
structural models. Fourth, I will prepare a theoretical paper or appendix that can support the
third-year empirical validation and eventual journal submission.

The main contribution of the second year is therefore not only theoretical. It provides the
identification logic that allows the full project to move from reduced-form uncertainty
measurement to structural economic interpretation.

\section{The Third Year: Methods, Procedures, and Implementation Schedule}

\subsection{Purpose of the Third Year}

The third year empirically validates the structural mechanisms derived in the second year using
Taiwan's aggregate data. The first-year OI-SVMVAR identifies external uncertainty transmission
and extracts Taiwan's macroeconomic and financial uncertainty factors. The second-year
wedge-projection theory derives the BCA signatures implied by competing mechanisms. The
third year connects these two components by testing whether Taiwan's empirical wedges respond
to uncertainty shocks in the way predicted by the theory.

This step is necessary because reduced-form evidence alone cannot distinguish structural
mechanisms. If financial uncertainty raises asset-market volatility and depresses investment, the
response may reflect a collateral constraint, a risk-premium channel, a working-capital channel,
or a combination of frictions. The BCA framework provides a way to separate these mechanisms
because different frictions project into different wedges. The third year therefore asks whether
Taiwan's observed wedge responses are closer to the investment-wedge signature of the collateral
or risk-premium channel, or to the labor-wedge signature of the working-capital channel.

\subsection{Constructing Taiwan's BCA Wedges}

The first task in the third year is to construct a small-open-economy BCA prototype for Taiwan.
The prototype decomposes Taiwan's aggregate fluctuations into four accounting wedges: an
efficiency wedge, a labor wedge, an investment wedge, and a government or external wedge. These
wedges summarize the distortions that a representative-agent prototype must introduce in order
to reproduce the observed paths of output, consumption, investment, labor input, and external
absorption.

The empirical construction will use Taiwan aggregate data from national accounts, labor-market
statistics, price indices, and external-sector accounts. The baseline BCA exercise will be
conducted at quarterly frequency because Taiwan's core national-account objects, including GDP,
private consumption, fixed capital formation, and external absorption, are measured most
reliably at that frequency. The monthly uncertainty factors from the first-year OI-SVMVAR will
therefore be time-aggregated to quarterly frequency before entering the wedge-level local
projections. The baseline aggregation will use quarterly averages of the monthly posterior draws,
with end-of-quarter values used as a robustness check for shocks that are more closely tied to
financial-market conditions. This design preserves the first-year model's ability to detect
short-lived uncertainty episodes while keeping the third-year BCA wedges anchored in reliable
national-account data.

The main empirical focus is on two wedges:
\[
\hat{\tau}_{x,t}
\qquad \text{and} \qquad
\hat{\Lambda}_{\ell,t}.
\]
The first object is the investment wedge. It captures distortions in the intertemporal Euler
equation and is the key empirical object for collateral and risk-premium mechanisms. The second
object is the after-tax labor wedge. It captures distortions in the intratemporal labor condition
and is the key empirical object for working-capital mechanisms.

The remaining BCA wedges are also informative. The efficiency wedge helps distinguish uncertainty
effects from productivity-like movements. The government or external wedge helps absorb changes
in public spending, trade balances, and external absorption. These wedges will be reported to
ensure that the mechanism comparison is not driven by omitted accounting components.

\subsection{Linking Uncertainty Factors to Wedge Signatures}

The core third-year exercise estimates how Taiwan's BCA wedges respond to the uncertainty
factors extracted in the first year. Let $w_t$ denote an empirical BCA wedge and let
$h_{m,t}$ and $h_{f,t}$ denote Taiwan's macroeconomic and financial uncertainty factors from
the OI-SVMVAR. The baseline local projection takes the form
\[
w_{t+h}
=
\alpha_h
+
\beta_{m,h} \varepsilon^h_{m,t}
+
\beta_{f,h} \varepsilon^h_{f,t}
+
\Gamma_h' X_t
+
u_{t+h},
\qquad h=0,1,\ldots,H,
\]
where $\varepsilon^h_{m,t}$ and $\varepsilon^h_{f,t}$ are macroeconomic and financial
uncertainty shocks extracted from the first-year model, and $X_t$ contains controls such as lagged
wedges, lagged macroeconomic variables, and external controls, while $u_{t+h}$ is the projection
residual. The coefficients
$\beta_{m,h}$ and $\beta_{f,h}$ trace the impulse responses of each wedge to macroeconomic and
financial uncertainty shocks.

The theoretical predictions from the second year guide the interpretation. If financial uncertainty
mainly raises the investment wedge, the evidence supports a collateral or risk-premium mechanism.
If financial uncertainty mainly lowers the after-tax labor wedge, the evidence supports a
working-capital mechanism. If macroeconomic uncertainty affects the efficiency wedge or the
external wedge more strongly, the evidence may point to real-side trade, productivity, or external
absorption channels.
This broader wedge comparison is important for Taiwan. If working-capital frictions operate
through imported materials or intermediate inputs rather than only through wage payments, the
response may appear partly in the efficiency wedge or the external wedge. I will therefore report
the full wedge vector and interpret labor-wedge evidence as the benchmark wage-finance case,
while treating efficiency- and external-wedge responses as evidence for input-finance or
trade-finance variants of the working-capital channel.

This wedge-level approach is more informative than matching impulse responses of standard
macroeconomic variables alone. A decline in investment, for example, can arise from many
mechanisms. The investment wedge asks whether the decline reflects an intertemporal distortion.
The labor wedge asks whether the shock acts like a distortion in the labor-demand or labor-supply
condition. The efficiency and external wedges help assess whether the response instead resembles
productivity or external-demand movements.

\subsection{MCMC Uncertainty Propagation}

A common limitation of two-step empirical designs is generated-regressor bias. If the uncertainty
factor is estimated in the first step and then treated as observed in the second step, the second-step
confidence intervals may be too narrow. This issue is especially relevant here because
$h_{m,t}$ and $h_{f,t}$ are latent factors estimated from a large Bayesian model.

To address this problem, the third-year analysis will propagate first-year posterior uncertainty
into the wedge-level local projections. I will not rely only on a single posterior mean path of
$h_{m,t}$ and $h_{f,t}$. Instead, I will use retained posterior draws from the first-year
OI-SVMVAR. For each posterior draw, I will compute the associated uncertainty shocks and
estimate the wedge responses. The final credible intervals will integrate over both the uncertainty
in the first-year factor estimates and the uncertainty in the third-year local projections.

This procedure has two advantages. First, it keeps the empirical validation consistent with the
Bayesian structure of the first-year model. Second, it makes the mechanism comparison more
conservative. If the data strongly support one mechanism even after posterior uncertainty is
propagated, the conclusion will be more credible. If the posterior uncertainty makes the evidence
ambiguous, the project will report the relative strength of evidence rather than forcing a binary
classification.
The procedure is also computationally feasible. The expensive step is the first-year OI-SVMVAR
estimation. Once posterior draws of the uncertainty factors are retained, the third-year local
projections are sequential linear regressions, so repeating them across posterior draws is much
less demanding than re-estimating the high-dimensional Bayesian VAR.

\subsection{Hypothesis Testing and Policy Interpretation}

The third-year empirical tests are organized around the wedge signatures derived in the second
year. The main hypotheses are:
\[
H_0^{\mathrm{coll}}:
\quad
\text{financial uncertainty mainly raises the investment wedge},
\]
and
\[
H_1^{\mathrm{wc}}:
\quad
\text{financial uncertainty mainly lowers the after-tax labor wedge}.
\]
The first hypothesis corresponds to an intertemporal collateral or risk-premium mechanism. The
second corresponds to a working-capital mechanism. The empirical analysis will compare the
magnitude, sign, timing, and statistical uncertainty of the wedge responses.

The policy implications differ across mechanisms. If the collateral or risk-premium channel
dominates, external uncertainty primarily affects Taiwan by tightening intertemporal financing
conditions. In that case, policy attention should focus on credit-market stabilization, liquidity
provision, and financial-market resilience. If the working-capital channel dominates, external
uncertainty acts more directly through firms' short-term financing costs and labor demand. In
that case, policies that reduce liquidity pressure on firms and stabilize working-capital finance
may be more relevant. If real-side wedges dominate, the appropriate response may involve export
support, industrial adjustment, or supply-chain resilience.

The empirical goal is not to assign all episodes to a single permanent mechanism. Taiwan's
transmission channel may change over time. For example, the Global Financial Crisis may display
a stronger financial-market component, while the U.S.--China trade war may show stronger real
and external-demand components. The third-year analysis will therefore report both average
responses and episode-specific evidence when the data allow.

\subsection{Third-Year Work Plan and Anticipated Results}

The third year is organized into four stages.

\begin{center}
\begin{tabular}{p{0.18\textwidth}p{0.36\textwidth}p{0.34\textwidth}}
\toprule
Period & Main task & Expected output \\
\midrule
Months 1--3 &
Construct Taiwan BCA prototype and wedge series &
Efficiency, labor, investment, and government/external wedges \\
\addlinespace
Months 3--6 &
Link Year 1 uncertainty posterior draws to wedge data &
Matched uncertainty-wedge dataset and posterior shock measures \\
\addlinespace
Months 6--9 &
Estimate MCMC-integrated wedge-level local projections &
Impulse responses of BCA wedges to uncertainty shocks \\
\addlinespace
Months 9--12 &
Test structural mechanisms and write final manuscript &
Mechanism comparison, policy interpretation, and journal-paper draft \\
\bottomrule
\end{tabular}
\end{center}

By the end of the third year, I expect to complete four outputs. First, I will produce a Taiwan
BCA wedge dataset suitable for studying uncertainty transmission. Second, I will estimate
wedge-level impulse responses to Taiwan's macroeconomic and financial uncertainty shocks.
Third, I will conduct a formal mechanism comparison using the theoretical signatures derived in
the second year. Fourth, I will prepare a complete manuscript that integrates the empirical
OI-SVMVAR evidence, the wedge-projection theory, and the Taiwan BCA validation.

The third year completes the project's main contribution. It transforms the first-year uncertainty
measurement exercise and the second-year theoretical mapping into direct evidence on the
structural mechanism through which external uncertainty affects Taiwan.

\section{Integrated Research Project}

\subsection{Integration Across the Three Years}

This project is designed as a unified three-year research program rather than three separate
studies. The first year identifies where external uncertainty comes from and whether it enters
Taiwan mainly through macroeconomic or financial channels. The second year turns those
empirical channels into structural predictions by deriving wedge signatures for collateral,
risk-premium, and working-capital mechanisms. The third year then tests these signatures in
Taiwan's BCA wedges while propagating posterior uncertainty from the first-year OI-SVMVAR.

The integration matters for policy because each year removes one layer of ambiguity. A finding
that U.S. monetary uncertainty raises Taiwan's financial uncertainty is informative, but it does
not by itself say whether the relevant policy margin is liquidity, collateral, credit access, or
real-side adjustment. The wedge-projection and BCA validation stages translate the empirical
classification evidence into mechanism-level diagnoses, so that the final project can distinguish
which policy instruments are most relevant under different external uncertainty episodes.

\subsection{Expected Academic Contributions}

The project makes four academic contributions.

First, it provides a new empirical account of external uncertainty transmission to Taiwan. Existing
work often studies uncertainty using a single proxy, such as financial volatility or economic policy
uncertainty. This project instead estimates domestic uncertainty factors jointly with their
macroeconomic effects and external drivers. This design allows the project to compare U.S.,
Chinese, global, and trade-policy sources within one coherent framework.

Second, the project adapts the OI-SVMVAR framework to a small open economy. The key
innovation is to reinterpret the unclassified-variable mechanism as a transmission-channel
identifier. Instead of pre-assigning external variables to macroeconomic or financial blocks, the
model lets the data determine whether each external driver is associated with Taiwan's
macroeconomic or financial uncertainty at each point in time. This approach provides a template
for studying other small open economies exposed to multiple external shocks.

Third, the project develops a wedge-projection framework for interpreting uncertainty shocks.
The framework shows how second-moment uncertainty corrections in nonlinear structural models
can appear as first-order BCA wedges. This contribution helps bridge reduced-form uncertainty
measurement and micro-founded structural interpretation.

Fourth, the project provides an empirical mechanism test using Taiwan's BCA wedges. Instead
of comparing only the impulse responses of standard macroeconomic variables, the project tests
whether uncertainty shocks move the wedges predicted by theory. This wedge-level validation
offers a sharper way to distinguish collateral, risk-premium, and working-capital channels.

\subsection{Expected Policy Contributions}

The project also has direct policy relevance for Taiwan. External uncertainty is a recurring feature
of Taiwan's economic environment. U.S. monetary tightening, China-related demand shocks,
trade-policy disputes, global financial volatility, and geopolitical risk all affect Taiwan, but they
do not necessarily operate through the same channel. A policy response that is appropriate for a
real-side demand shock may be less effective for a financial-market shock.

By distinguishing macroeconomic and financial transmission channels, the project can help
policymakers interpret uncertainty episodes more precisely. If an external shock mainly affects
Taiwan through macroeconomic channels, the relevant policy tools are closer to MOEA and NDC
instruments, including export credit guarantees, small and medium enterprise (SME) liquidity
support, supply-chain resilience programs, and industrial adjustment measures. If the shock
mainly affects Taiwan through financial channels, the relevant tools are closer to CBC liquidity
operations, macroprudential measures, selective credit controls, exchange-market monitoring,
and Financial Supervisory Commission (FSC) monitoring of market functioning and
financial-stability risks.

The third-year wedge validation further improves the policy interpretation. If the evidence points
to a collateral or risk-premium channel, CBC and FSC financial-stability tools become more
central. If the evidence points to a working-capital channel, SME liquidity facilities, export credit
guarantees, and bank-credit support become more relevant. If real-side wedges dominate, MOEA
and NDC export promotion, supply-chain adjustment, and industrial upgrading policies may be
more appropriate.

The project therefore contributes not only to academic research on uncertainty shocks, but also
to the practical diagnosis of Taiwan's external vulnerability.

\subsection{Expected Outputs and Deliverables}

The expected outputs are organized by year.

\begin{center}
\begin{tabular}{p{0.15\textwidth}p{0.38\textwidth}p{0.35\textwidth}}
\toprule
Year & Main research output & Deliverables \\
\midrule
Year 1 &
Empirical identification of Taiwan's uncertainty transmission &
Monthly dataset, OI-SVMVAR estimates, uncertainty factors, classification
probabilities, IRFs, FEVDs, conference-paper draft \\
\addlinespace
Year 2 &
Theoretical wedge-projection framework &
Formal theorem, wedge-signature table, numerical validation, theoretical paper
or appendix \\
\addlinespace
Year 3 &
Empirical BCA validation and mechanism testing &
Taiwan BCA wedge dataset, MCMC-integrated local projections, mechanism
comparison, journal-paper draft \\
\bottomrule
\end{tabular}
\end{center}

The final project will produce three main research products. The first is a Taiwan external
uncertainty database and empirical paper. The second is a theoretical wedge-projection paper or
appendix. The third is an integrated empirical-structural paper on Taiwan's external uncertainty
transmission.

\subsection{Risk Management: Anticipated Problems and Means of Resolution}

The project is ambitious, but its risks are manageable. The main risks and mitigation strategies
are summarized below.

\begin{center}
\begin{tabular}{p{0.28\textwidth}p{0.58\textwidth}}
\toprule
Risk & Mitigation strategy \\
\midrule
The full 52-variable OI-SVMVAR is slow to converge. &
Begin with a reduced benchmark model using the core Taiwan macroeconomic and financial
anchors and the most important external drivers. Move to the full model after validating the
sampler and convergence diagnostics. \\
\addlinespace
Some data series require interpolation or have limited coverage. &
Report the baseline specification transparently, maintain a complete data-construction log, and
conduct robustness checks that exclude interpolated variables or use alternative sample windows. \\
\addlinespace
China-related monthly data are difficult to harmonize. &
Use the downloaded China industrial production, PPI, and M2 series after source validation, and
compare them with alternative public sources when available. If definitions or vintages remain
fragile, report a reduced China block as a robustness specification while keeping U.S. and global
indicators in the baseline. \\
\addlinespace
Theoretical wedge signatures are not sharply separated in all cases. &
Report the relative strength of evidence across mechanisms rather than forcing a binary
classification. Use the timing, sign, and wedge location of the responses jointly. \\
\addlinespace
Local projection results are sensitive to horizon choice. &
Report full horizon profiles, including impact responses, peak responses, and sign changes. Avoid
interpreting only the maximum response at a long horizon. \\
\addlinespace
First-stage uncertainty factors create generated-regressor concerns. &
Propagate posterior uncertainty from the first-year OI-SVMVAR into the third-year wedge-level
local projections instead of using only posterior mean paths. \\
\bottomrule
\end{tabular}
\end{center}

These safeguards ensure that the project remains feasible even if some components require
adjustment. The reduced benchmark model protects the first-year empirical analysis. The
theoretical wedge-signature table protects the second-year mechanism comparison. The
MCMC-integrated validation protects the third-year inference. Together, these safeguards make
the project suitable for a three-year research grant.

\subsection{Concluding Statement}

This project studies Taiwan as a small open economy caught between major external sources of
uncertainty. Its contribution is to move beyond asking whether uncertainty matters. Instead, it
asks where uncertainty comes from, how it enters Taiwan, and which structural friction explains
its effects.

The first year provides empirical identification through an order-invariant large-scale
SVMVAR. The second year provides theoretical interpretation through wedge projection. The
third year provides empirical validation through Taiwan's BCA wedges. By combining these three
components, the project delivers a coherent framework for measuring, interpreting, and validating
external uncertainty transmission in Taiwan.

\appendix
\section{Econometric Model: OI-SVMVAR}
\label{app:oi-svmvar}

This appendix summarizes the order-invariant stochastic volatility in mean vector autoregression
(OI-SVMVAR) used in the first-year empirical analysis. Let $\mathbf{y}_t$ denote the $N \times 1$
vector of transformed and standardized monthly observables, partitioned into macroeconomic,
unclassified, and financial blocks. The observation equation is
\begin{equation}
\mathbf{B}_0 \mathbf{y}_t
= \sum_{\ell=1}^{p} \mathbf{B}_{\ell}\mathbf{y}_{t-\ell}
+ \sum_{j=0}^{q} \mathbf{A}^{h}_{j}\mathbf{h}_{t-j}
+ \boldsymbol{\varepsilon}_t,\qquad
\boldsymbol{\varepsilon}_t \sim N(\mathbf{0},\mathbf{U}_t),
\label{app:obs}
\end{equation}
where $\mathbf{h}_t=(h_{m,t},h_{f,t})'$ contains Taiwan's common macroeconomic and
financial log-volatility factors. The matrices $\mathbf{A}^{h}_{j}$ capture the stochastic-volatility-
in-mean channel: changes in uncertainty can directly shift the conditional mean of output,
employment, asset prices, credit conditions, and other observables. The contemporaneous impact
matrix $\mathbf{B}_0$ is unrestricted apart from the normalization of its diagonal elements to one.
This is the key order-invariance feature of the model: inference is not tied to a Cholesky ordering
of the variables in $\mathbf{y}_t$.

The conditional variance matrix in the structural equation is diagonal and block-partitioned as
\begin{equation}
\mathbf{U}_t =
\begin{pmatrix}
\mathbf{U}_{m,t} & \mathbf{0} & \mathbf{0} \\
\mathbf{0} & \mathbf{U}_{u,t} & \mathbf{0} \\
\mathbf{0} & \mathbf{0} & \mathbf{U}_{f,t}
\end{pmatrix},\qquad
\mathbf{U}_{k,t}=\operatorname{diag}\!\left(e^{\omega^k_{1,t}},\ldots,
e^{\omega^k_{n_k,t}}\right),\quad k\in\{m,u,f\}.
\label{app:ut}
\end{equation}
The block ordering is not an identifying recursive ordering; it only records whether a variable is
anchored as macroeconomic, left unclassified, or anchored as financial. Contemporaneous
interactions among all variables remain unrestricted through $\mathbf{B}_0$.

The log-volatility of each observable is decomposed into an idiosyncratic component and one of
the two common uncertainty factors:
\begin{align}
\omega^m_{i,t} &= \eta^m_{i,t}+h_{m,t}, \qquad i=1,\ldots,n_m, \nonumber\\
\omega^u_{i,t} &= \eta^u_{i,t}+h_{s_{i,t},t}, \qquad s_{i,t}\in\{m,f\},\quad i=1,\ldots,n_u,
\label{app:logvol}\\
\omega^f_{i,t} &= \eta^f_{i,t}+h_{f,t}, \qquad i=1,\ldots,n_f. \nonumber
\end{align}
The idiosyncratic terms $\eta^k_{i,t}$ capture variable-specific volatility movements. Variables
classified ex ante as macroeconomic load on $h_{m,t}$, and variables classified ex ante as
financial load on $h_{f,t}$. For unclassified variables, including external U.S., Chinese, global,
and policy-risk indicators, the latent state $s_{i,t}$ determines whether the variable's common
volatility component is assigned to the macroeconomic or financial factor at time $t$.

The two common factors evolve jointly according to
\begin{equation}
\mathbf{h}_t
= \sum_{r=1}^{p_h} \boldsymbol{\Phi}_r \mathbf{h}_{t-r}
+ \sum_{s=1}^{p_y} \boldsymbol{\Psi}_s \mathbf{y}_{t-s}
+ \boldsymbol{\varepsilon}^{h}_t,\qquad
\boldsymbol{\varepsilon}^{h}_t \sim N(\mathbf{0},\boldsymbol{\Sigma}_h).
\label{app:hdyn}
\end{equation}
This specification allows macroeconomic and financial uncertainty to interact over time and
allows lagged domestic and external conditions to forecast Taiwan's uncertainty factors. The
factors are therefore estimated jointly with the VAR rather than constructed in a preliminary
step.

For each unclassified variable $i$, the state $s_{i,t}$ follows a two-state first-order Markov chain.
The main transmission-channel object is the posterior classification probability
\begin{equation}
\pi_{i,t}
=P(s_{i,t}=m\mid \mathbf{y}_{1:T}),
\label{app:pi}
\end{equation}
which is the posterior probability that variable $i$ belongs to the macroeconomic block at time
$t$. Values of $\pi_{i,t}$ near one indicate macroeconomic transmission, while values near zero
indicate financial transmission. Because $\pi_{i,t}$ is time-varying, the same external driver can
operate through different channels in different episodes.

The model is estimated by Bayesian MCMC following the OI-SVMVAR algorithm of
\citet{davidson2025investigating}. The sampler cycles through the VAR and volatility-in-mean
coefficients, the unrestricted contemporaneous matrix $\mathbf{B}_0$, the common and
idiosyncratic log-volatility paths, and the latent classification paths. The unrestricted
$\mathbf{B}_0$ is sampled using the order-invariant row-wise parameter transformation developed
for the model, preserving permutation-invariant inference without imposing block exogeneity or
recursive zero restrictions. The finite-state Markov paths $(s_{i,1},\ldots,s_{i,T})$ and their
transition probabilities are drawn with forward-filtering backward-sampling and Chib-style path
sampling. Posterior draws of $\mathbf{h}_t$ and $\pi_{i,t}$ are then used for the classification,
variance-decomposition, and impulse-response analyses.

\section{Economic Model: Structural Environment and Wedge Projection}
\label{app:wedge-projection}

This appendix records the compact economic model behind the wedge-projection exercise. The
purpose is not to estimate a full DSGE model inside the proposal, but to make clear which
structural environment generates the wedge signatures used in the second and third years.
The baseline is a Fern\'andez-Villaverde--Guerr\'on-Quintana-style two-agent collateral economy
with stochastic volatility in real, demand, and financial-friction shocks. I then contrast this
intertemporal channel with a working-capital alternative.

Time is discrete. A household contains workers and entrepreneurs. The household head allocates
equity claims and writes contingent plans before idiosyncratic member types are realized.
Workers supply labor and receive wages; entrepreneurs invest in capital and face a collateral
constraint. A compact representation of preferences is
\begin{equation}
\E_0\sum_{t=0}^{\infty}\beta^t d_t
\left[
\pi \frac{(c^E_t)^{1-\rho}-1}{1-\rho}
+(1-\pi)\frac{\{c^W_t(1-\ell_t)^\eta\}^{1-\rho}-1}{1-\rho}
\right],
\label{app:fvgq-utility}
\end{equation}
where $c^E_t$ and $c^W_t$ are entrepreneur and worker consumption, $\ell_t$ is labor,
$d_t$ is a preference shifter, and $\pi$ is the fraction of entrepreneurs.

Final output is produced competitively using utilized capital and labor:
\begin{equation}
y_t=A_t(u_t k_t)^\alpha \ell_t^{1-\alpha},
\qquad
r_t=\alpha\frac{y_t}{u_t k_t},
\qquad
w_t=(1-\alpha)\frac{y_t}{\ell_t}.
\label{app:fvgq-production}
\end{equation}
Capital utilization $u_t$ raises depreciation through a convex schedule
\begin{equation}
\delta(u_t)=\delta_0+\delta_1(u_t-1)+\frac{\delta_2}{2}(u_t-1)^2,
\qquad
k_{t+1}=(1-\delta(u_t))k_t+\pi i_t.
\label{app:fvgq-capital}
\end{equation}
The key financial friction is a skin-in-the-game collateral constraint. If $s^E_{t+1}$ denotes
entrepreneur equity claims carried into the next period, then
\begin{equation}
s^E_{t+1}
\ge
(1-\theta_t)i_t+(1-\phi)(1-\delta(u_t))s_t.
\label{app:fvgq-collateral}
\end{equation}
The time-varying pledgeability parameter $\theta_t$ controls how much new capital can be
externally financed, while $\phi$ controls resaleability of existing capital. When the constraint
binds, changes in uncertainty alter the intertemporal Euler equation through collateral-risk,
risk-premium, and precautionary terms.

The structural shocks are a real productivity shock $A_t$, a demand shifter $d_t$, and a financial
friction shock $\theta_t$. Each follows a persistent level process with stochastic volatility:
\begin{align}
z_{j,t} &=(1-\rho_j)\bar z_j+\rho_j z_{j,t-1}
       +\exp(\sigma_{j,t})\varepsilon_{j,t},
       \qquad \varepsilon_{j,t}\sim N(0,1), \nonumber\\
\sigma_{j,t} &=(1-\rho_{\sigma_j})\bar\sigma_j
       +\rho_{\sigma_j}\sigma_{j,t-1}
       +(1-\rho_{\sigma_j}^2)^{1/2}\nu_j u_{j,t},
       \qquad u_{j,t}\sim N(0,1),
\label{app:fvgq-sv}
\end{align}
for $j\in\{A,d,\theta\}$. The volatility innovations $u_{j,t}$ are the structural uncertainty
shocks. In the empirical project, the OI-SVMVAR uncertainty factors provide observed
counterparts to these latent volatility states.

In a representative-household BCA prototype, the log-linearized investment Euler equation can
be written as
\begin{equation}
\hat{\tau}_{x,t}-\gamma \tilde C_t
=
\E_t\left[
-\gamma \tilde C_{t+1}
+\omega_K \tilde R^k_{t+1}
+\omega_\tau \hat{\tau}_{x,t+1}
\right],
\label{app:bca-euler}
\end{equation}
where $\hat{\tau}_{x,t}$ is the investment wedge, $\tilde C_t$ is measured
consumption, $\tilde R^k_{t+1}$ is the measured capital return, $\gamma$ is the risk-aversion
parameter, and $\omega_K,\omega_\tau$ are steady-state coefficients. If the true data-generating process
contains stochastic volatility, the corresponding Euler equation has a second-order risk
correction:
\begin{equation}
-\gamma \tilde C_t
=
\E_t\left[
-\gamma \tilde C_{t+1}
+\omega_K \tilde R^k_{t+1}
\right]
+\frac{1}{2}\Phi(\Omega_t)
+o(\|\Delta\Omega_t\|),
\label{app:true-euler}
\end{equation}
where $\Omega_t$ denotes the volatility state and $\Phi(\Omega_t)$ collects conditional
variance, covariance, and Hessian terms. Comparing \eqref{app:bca-euler} and
\eqref{app:true-euler} yields the wedge projection equation
\begin{equation}
\hat{\tau}_{x,t}
=
\omega_\tau \E_t \hat{\tau}_{x,t+1}
+\frac{1}{2}\Phi(\Omega_t).
\label{app:projection}
\end{equation}

Let $\Delta\Omega_t=\Omega_t-\bar{\Omega}$ and suppose the volatility state follows a stable
matrix AR(1),
\begin{equation}
\Delta\Omega_{t+1}=R_\Omega\Delta\Omega_t+\eta_{t+1},
\qquad \rho(R_\Omega)<1,
\label{app:matrix-ar}
\end{equation}
with local approximation
$\Phi(\Omega_t)=\Phi(\bar{\Omega})+\bm{\phi}'\Delta\Omega_t+
o(\|\Delta\Omega_t\|)$. If $|\omega_\tau|<1$, the locally bounded solution is
\begin{equation}
\hat{\tau}_{x,t}
=
\frac{\Phi(\bar{\Omega})}{2(1-\omega_\tau)}
+\frac{1}{2}\bm{\phi}'(I-\omega_\tau R_\Omega)^{-1}\Delta\Omega_t
+o(\|\Delta\Omega_t\|).
\label{app:closed-form}
\end{equation}
For a scalar AR(1), $\Delta\Omega_{t+1}=\rho_\Omega\Delta\Omega_t+\eta_{t+1}$, this becomes
\begin{equation}
\hat{\tau}_{x,t}
=
\frac{\Phi(\bar{\Omega})}{2(1-\omega_\tau)}
+\underbrace{\frac{\phi}{2(1-\omega_\tau\rho_\Omega)}}_{\kappa_x}
(\Omega_t-\bar{\Omega})
+o(|\Omega_t-\bar{\Omega}|).
\label{app:scalar-closed-form}
\end{equation}
Thus a purely second-moment structural force may be recovered by BCA as a persistent
first-order investment wedge.

In this collateral model, the risk correction can be organized schematically as
\begin{equation}
\Phi^{collat}(\Omega_t)
=
\Phi^{prec}(\Omega_t)
+\Phi^{rp}(\Omega_t)
+\Phi^{coll}(\Omega_t)
+\Phi^{agg}(\Omega_t).
\label{app:fvgq-decomposition}
\end{equation}
The precautionary term is increasing in conditional consumption risk,
$\Phi^{prec}_t\propto \gamma^2 \mathrm{Var}_t(\hat c_{t+1})$. The risk-premium term depends
on the consumption-return covariance,
$\Phi^{rp}_t\propto -\gamma\,\mathrm{Cov}_t(\hat c_{t+1},\hat r^k_{t+1})$. The collateral
term captures the effect of volatility on pledgeability or collateral tightness,
$\Phi^{coll}_t\propto \nu_\theta\sigma_{\theta,t}^2$, where $\nu_\theta$ governs the sensitivity
of the collateral constraint to financial-friction volatility. In a two-agent version, aggregation
adds a dispersion term, for example
$\Phi^{agg}_t\propto \mathrm{Var}_t(\hat c_{E,t+1}-\hat c_{W,t+1})$, because the Euler
equation pricing investment is tied to entrepreneur marginal utility while BCA uses aggregate
consumption. These components imply that the sign of $\kappa_x$ is disciplined by the
relative strength of precautionary motives, risk premia, collateral risk, and aggregation.

A working-capital mechanism has a different signature. Suppose firms finance a fraction
$\theta_{wc}$ of the wage bill in advance and face external finance spread $q_t$. The true
static labor condition is
\begin{equation}
A_tF_{L,t}=w_t(1+\theta_{wc}q_t),
\label{app:wc-foc}
\end{equation}
where $F_{L,t}$ is the marginal product of labor. The BCA intratemporal condition is
\begin{equation}
-\frac{U_L(C_t,L_t)}{U_C(C_t,L_t)}
=
\Lambda_{\ell,t}^{BCA} A_tF_{L,t}.
\label{app:bca-labor}
\end{equation}
Here $U_C$ and $U_L$ denote the marginal utilities of consumption and labor.
Using the household condition $-U_L/U_C=w_t$ gives the exact projection
\begin{equation}
\Lambda_{\ell,t}^{BCA}
=
\frac{1}{1+\theta_{wc}q_t}.
\label{app:wc-projection}
\end{equation}
Around $\bar q$,
\begin{equation}
\hat{\Lambda}_{\ell,t}^{BCA}
=
-\frac{\theta_{wc}}{1+\theta_{wc}\bar q}(q_t-\bar q)
+O((q_t-\bar q)^2).
\label{app:wc-linear}
\end{equation}
If $q_t-\bar q=\alpha_\xi\hat\xi_t+\alpha_\Omega(\Omega_t-\bar{\Omega})+o(\cdot)$ with
$\alpha_\xi>0$ and $\alpha_\Omega\ge0$, a financial tightening or higher uncertainty lowers
the after-tax labor wedge. The model therefore distinguishes an intertemporal
collateral/risk-premium channel, which projects mainly into $\hat{\tau}_{x,t}$, from a
working-capital channel, which projects mainly into $\hat{\Lambda}_{\ell,t}^{BCA}$.

\bibliographystyle{aea}
\bibliography{references}

\end{document}
